\documentclass[output=paper]{langscibook}
\ChapterDOI{10.5281/zenodo.22880240}
\title{\textit{O sinivi} – The canoe tree (Blackboard Tree)}
\author{Joyce Maion}
\abstract{}
\begin{document}
\maketitle\label{chap:7}
\begin{figure}[H]
    \includegraphics[width=.40\linewidth]{figures/7_1.jpg}\hfil
    \includegraphics[width=.33\linewidth]{figures/7_2.jpg}
    \caption{\textit{O naono sinivi} -- Canoe tree}
    \label{fig:7_1-2}
\end{figure}
\begin{figure}[H]
    \includegraphics[width=.6\linewidth]{figures/7_3.jpg}
    \caption{\textit{Paka sinivi} -- Leaf of canoe tree}
    \label{fig:7_3}
\end{figure}
\begin{xltabular}{\textwidth}{lXX}
    (\ref{ex:text3-1})&\textit{\textbf{O vahutate no naono sinivi.}}&\textbf{The story of the canoe tree (Blackboard Tree.)}\\
    (\ref{ex:text3-2})&\textit{O sinivi na paku raara bona maa naono to dao raara bono maramakasi, o bio ae o naosinivi.}&We make the canoe from the trees that are called \textit{maramakasi}, \textit{bio} (a Ficus species), and canoe tree.\footnote{A further tree used for canoes and mentioned in the TLC is a terminalia species called \textit{marakaoto}, see \url{https://dictionaria.clld.org/units/teop-marakaoto}}\\
    (\ref{ex:text3-3})&\textit{Evehee enaa kahi vahutate koa ni nom o naosinivi.}&But I'll talk only about the canoe tree.\\
    (\ref{ex:text3-4})&\textit{O naono vai o sinivi na taba beera kurusu kiri nana bona tavaan to popo rori teo toro mohina.}&This canoe tree is a very important thing for the people who live on an island.\\
    (\ref{ex:text3-5})&\textit{Bea tavaan tara bono naono sinivi o sii rutaa beve pura nana rasuu,}&When the people see a little canoe tree growing in the jungle,\\
    (\ref{ex:text3-6})&\textit{eori repaa kiu kahi ni bari,}&they clear the ground around it,\\
    (\ref{ex:text3-7})&\textit{ore paa sun nana tea mamana.}&so that it stands in the open.\\
    (\ref{ex:text3-8})&\textit{Eve repaa pura hiava komana, beori kikira batari bari.}&Then it grows up, while they are looking after it.\\
    (\ref{ex:text3-9})&\textit{Beo meho aba tara bono naosinivi paa kiu kahi ni vori,}&If other people see that the canoe tree has been freed from the undergrowth around it,\\
    (\ref{ex:text3-10})&\textit{ori paa kona voen, a aba to paa kiu kahi niu bona, eve to vaahito nana. Saka antee haa riori tea mosi.}&they recognize that the person who has done this work is the one who owns it. They cannot cut it down.\\
    (\ref{ex:text3-11})&\textit{Eori to popo rori rapana kasuana bara tea maa toro mohina, toro tei mi bono sinivi,}&Those who live along the coast or on the islands must have a canoe,\\
    (\ref{ex:text3-12})&\textit{tea tabae beori hiki, eori repaa hiki ta taba toro hae ori tea vagana.}&because if they don't, they wouldn't have anything to board and go fishing.\\
    (\ref{ex:text3-13})&\textit{Beori rake mee tea hua pahana tea maa meha maa toro mohina.}&Also, when they want to paddle across the bay to the islands.\\
    (\ref{ex:text3-14})&\textit{Bea aba hiki to naosinivi mene paku sinivi,}&If someone does not have a canoe tree for making a canoe,\\
    (\ref{ex:text3-15})&\textit{toro paku ta hanana vai toro vaahito kaakoo me to sinivi kiri bona vuu tom sinana.}&he or she should find some other way to acquire a canoe for the family.\\
    (\ref{ex:text3-16})&\textit{Meara toro kikira nio naosinivi eara mene hiki sinivi.}&We must look after the canoe tree, otherwise we would be without canoes.\\
\end{xltabular}
\begin{figure}[H]
    \includegraphics[width=.9\linewidth]{figures/7_4.jpg}
    \caption{\textit{Sinivi suuna} -- Indigenous canoe}
    \label{fig:7_4}
\end{figure}
\end{document}